\section{导读：上市不是终点，而是一次路线复盘}

这期访谈发生在一个特殊节点：智谱确认于 2026 年 1 月 8 日登陆港交所，成为中国首家上市的大模型公司，也被称为“全球大模型第一股”。张鹏带着摔伤的腿、拄着拐杖接受访谈，节目用 “Break a leg” 做开场隐喻。但这期真正重要的，不是上市敲钟的戏剧性，而是智谱六年多如何从清华实验室的认知智能探索，走到大模型产业和上市公司治理。

这份笔记把它整理成三条线。第一条是组织线：清华知识工程实验室的 P2P 传统，怎样变成智谱的成果转化路径。第二条是技术线：从感知智能到认知智能，再被 GPT-3、ChatGPT、DeepSeek 和 Scaling Law 反复冲击。第三条是公司线：一家学院派 AI 公司如何处理开源与闭源、理想主义与商业化、AGI 长期目标与上市公司短期纪律。

\begin{importantbox}{本期核心命题}
智谱的故事不是“突然赶上大模型风口”，而是长期把研究、工程、产业和商业化放在同一条链路里。上市只是阶段性结果；张鹏更希望智谱在历史里被写成“AGI 的先行者”和“开路的人”。
\end{importantbox}

\begin{knowledgebox}{视觉策略说明}
本视频是固定访谈画面，没有 slides、白板或产品演示。正文不重复插入人物帧；封面用于来源识别，正文用成果转化、模型冲击、Scaling Law、开源闭源和公司阶段等概念图承载教学内容。
\end{knowledgebox}

\subsection{本章小结}

EP129 适合和 EP136 的大模型季报、EP139 的 Agent 技术史一起读。它从一家模型公司的内部视角，解释大模型产业如何被技术范式、资本市场、开源生态和组织治理共同塑形。

